\section{The Law Written Within}

The command has reached the will without yet healing it. Scripture's next move is not to relax the good but to promise a heart capable of embracing it.

The prophets promise something more than a clearer inscription. Jeremiah speaks of a covenant in which God's law is written upon the heart. Ezekiel joins a new heart to the gift of God's Spirit, who enables the people to walk in the divine ways (Jer 31:31--34; Ezek 36:26--27). The promised transformation neither declares the commandments mistaken nor treats inward sincerity as a substitute for them. It brings the will into agreement with the good they name.

Augustine makes this the center of \emph{On the Spirit and the Letter}. The law without assisting grace stands outside the will: it teaches and accuses, but cannot justify. When the life-giving Spirit is present, what was feared as an external demand becomes loved as an inward law. Augustine identifies the law written by God in the heart with the presence of the Holy Spirit, through whom charity is poured into us and the law is fulfilled (chs.~32--36). His famous petition in the \emph{Confessions} is not an escape from responsibility: ``Give what you command, and command what you will'' (X.29.40). The command remains entire. Its fulfillment is confessed as gift.

Aquinas gathers that Augustinian claim into a definition of the New Law. What is chief in the law of the Gospel is not a text but the grace of the Holy Spirit given through faith in Christ. Written teaching remains necessary: it disposes persons to receive grace and directs them in its use. Yet the New Law does more than indicate what should be done. By grace it helps the person to do it (ST I--II, q.~106, a.~1).

Charity is therefore neither an achievement assembled from repeated natural acts nor a divine feeling substituted for human willing. It is infused by the Holy Spirit because the act by which God is loved above all things for God's own sake exceeds the proportion of created nature (ST II--II, q.~24, a.~2). Grace heals the will wounded by sin and elevates it to a supernatural end.

The word \emph{infused} can sound as though the person were merely a vessel. But God moves the will according to its voluntary mode: even the gratuitous inward motion that prepares a person for grace gives rise to a genuinely voluntary act (ST I--II, q.~10, a.~4; q.~109, a.~6). Grace does not love while the person remains inert. It enables a free act that nature cannot manufacture.

The law written within is not a private oracle. It is the grace of Christ communicated through the Holy Spirit, received in faith and the sacramental life, and ordered toward acts that can be tested against the revealed command. An inward impulse that contradicts truth, justice, or love of neighbor does not become charity by feeling spontaneous.

Grace does not silence the command. It enables willing obedience from within. What relation to God can such obedience become?
